\section{三重積の対数型の非可換双有理変換}
\label{sec:cube-log}

三重積の三つの対数型成分の一般元は、因子を並べ替えた後、
\[
 \pi=(a x^2+b xz+c z^2)\partial_x\wedge\partial_z
       +\kappa y\partial_y\wedge(x\partial_x+z\partial_z)
\]
と正規化される。ここで $b^2-4ac\ne0$ を仮定する。
同梱の量子化原稿第4節では同変な cubic $\mathbb Z$-代数を用いた。
この節では分数体が直接計算できる別の斉次表示を構成する。

$B_R=R[Y,Z]$ とし、$R$ は必要な非零多項式を逆にした
$\mathbb C[\hbar]$ の局所化とする。
\[
 L(Y)=\tfrac b2Y+cZ,\qquad
 L(Z)=-aY-\tfrac b2Z,\qquad
 \sigma=1+\hbar L\quad\text{on }\langle Y,Z\rangle.
\]
$\det\sigma$ を逆にして $\sigma$ を $B_R$ の自己同型にし、$t=YZ$ と置く。
generalized Weyl algebra
\[
 G_R=B_R\langle X,W\rangle/
 \bigl(Xf-\sigma(f)X,\;Wf-\sigma^{-1}(f)W,\;
       WX-t,\;XW-\sigma(t)\bigr)
\]
の全生成元に次数1を与える。

\begin{proposition}\label{prop:gwa-quadric}
$G_R$ は次数ごとに自由な noetherian domain で、Hilbert 級数は
$(1+s)/(1-s)^3$ である。$\hbar=0$ では
$\mathbb C[X,Y,Z,W]/(XW-YZ)$ になる。
基底を $0$ の近傍に制限すると、各幾何学的ファイバーの射影圏は
smooth な非可換 quadric であり、その半古典極限は
$(a x^2+b xz+c z^2)\partial_x\wedge\partial_z$ である。
\end{proposition}
\begin{proof}
$G_R$ は $B_R[X^{\pm1};\sigma]$ に $W=tX^{-1}$ として埋め込まれ、
\[
 B_R\ \oplus\!\bigoplus_{n\ge1} B_RX^n
       \ \oplus\!\bigoplus_{n\ge1} B_RW^n
\]
が正規形となる。これが平坦性、domain 性、Hilbert 級数を与える。
さらに
\[
 S_R=B_R[X;\sigma][W;\tau,\delta],\quad
 \tau|_{B_R}=\sigma^{-1},\quad\tau(X)=X,\quad
 \delta|_{B_R}=0,\quad\delta(X)=t-\sigma(t)
\]
と置くと、$\tau,\delta$ は最初の Ore 関係を保存する。
$g=WX-t$ は中心的な正則二次元で、$G_R=S_R/(g)$ である。
各ファイバーの $S_R$ は四次元の homogeneous Ore AS-regular 代数である。
これは Ore 拡大の twisted Calabi--Yau 性と正の次数付けからも従う
\cite[定理0.2]{LiuWangWu}。
$S_b$ は $Y,Z,X,W$ の順序に関する quadratic PBW 代数なので Koszul であり、
中心的な正則二次元による商 $G_b=S_b/(g)$ も Koszul である。

smoothness を有限次元の条件で確認する。
Smith--Van den Bergh \cite[補題5.1、命題5.2、定理5.6]{SmithVdB}
により、$G_R$ の quadratic dual に付随する有限次元代数
\[
 C(G_R)=G_R^![w^{-1}]_0
\]
が半単純なら $\operatorname{qgr}G_R$ は smooth である。
ここで $w$ は $G_R^!$ の中心的な正則次数2の元であり、
$G_R^!/(w)=S_R^!$ である。
各ファイバーで $\dim(G_R^!)_4=\dim(G_R^!)_8=8$ で、
$w^2:(G_R^!)_4\to(G_R^!)_8$ は同型である。
従って基底を縮めれば、$C(G_R)$ は $(G_R^!)_4$ 上にこの同型で
積を戻した階数8の代数族として構成できる。
古典点では通常の smooth quadric の even Clifford 代数
$M_2(\mathbb C)\oplus M_2(\mathbb C)$ である。
正則表現の trace pairing の行列式が非零である開集合を取れば、
全ファイバーで半単純性と smoothness が従う。

半古典括弧は
$\{Y,Z\}=0$、$\{X,f\}=L(f)X$、$\{W,f\}=-L(f)W$ である。
$x=Y/X,z=Z/X$ とすると
\[
 \{x,z\}=z\bigl(\tfrac b2x+cz\bigr)
             -x\bigl(-ax-\tfrac b2z\bigr)
          =ax^2+bxz+cz^2.
\]
\end{proof}

重み $(X,Y,Z,W)=(0,1,1,2)$ の作用 $\rho$ は上の関係を保つ。
$p=1+\hbar\kappa$ をさらに $R$ 上で可逆にして、
\[
 H_R=G_R[Y_0][Y_1;\rho_p],\qquad
 \mathcal R_R=\bigoplus_{n\ge0}(H_R)_{n,n}
\]
と置く。$Y_0$ は中心的であり、$Y_i$ は第二次数1を持つ。
$\mathcal R_R$ は階数 $(n+1)^3$ の次数ごとに平坦な斉次環で、
古典点は $(\PP^1)^3$ の Segre 環である。
同変 Ore 拡大と Segre 積には補題\ref{lem:ore-projective-bundle}を適用できる。
各幾何学的ファイバー $G_b$ の dg smoothness と properness は、Orlov の分解
\cite[定理2.5]{OrlovSing} と Smith--Van den Bergh の
MCM 圏と $D^b(\operatorname{mod}C(G_b))$ の同値から従う。
実際、Gorenstein parameter は2なので二つの構造層 twist と
$C(G_b)$ の二つの単純加群が full exceptional collection を与える。
従って得られる三次元模型は次数ごとに平坦で、
各幾何学的ファイバーが dg smooth proper である。
アフィン座標 $y=Y_1/Y_0$ の交換子は混合項
$\kappa y\partial_y\wedge(x\partial_x+z\partial_z)$ を与える。

\begin{theorem}\label{thm:cube-log-torus}
上の射影量子化の division ring は、
一つの交換係数が1である三次元量子トーラスの division ring と同型である。
従って三重積の対数型の一般元には、トーラス型と非可換双有理な
射影量子化を選べる。
\end{theorem}
\begin{proof}
$L$ の固有一次式を $U,V$、固有値を $\nu,-\nu$ とする。
$\nu^2=(b^2-4ac)/4$ である。
$\sigma(U)=\lambda U,\sigma(V)=\mu V$、
$\lambda=1+\hbar\nu,\mu=1-\hbar\nu$ と置く。
次数零の分数体は $r=UX^{-1},s=VX^{-1}$ で生成され、
\[
 rs=\frac\lambda\mu sr,
 \qquad yr=p\,ry,\qquad ys=p\,sy.
\]
実際、$W=YZ X^{-1}$ なので $W$ は追加の分数体生成元を与えない。
$a_0=rs^{-1}$ と置けば
\[
 ya_0=a_0y,\qquad
 a_0s=\frac\lambda\mu sa_0,\qquad ys=p\,sy.
\]
これが所望の量子トーラス表示である。
命題\ref{prop:toric-common-field}により、同じ分数体を持つ
三模型のトーラス型 compactification が得られる。
\end{proof}

パラメーターを $q=\mu/\lambda$ と合わせると、
この分数体は第\ref{sec:line-conic-quantum}節の line--conic 型とも一致する。
この対応では量子化パラメーターの高次項も上式によって選んでいる。

